%==============================================================================
% Certificats de primalité inconditionnels pour p = 3m(m+1)+1
% H. Bakkaoui — 2026
%==============================================================================
\documentclass[11pt,a4paper]{article}

\usepackage[utf8]{inputenc}
\usepackage[T1]{fontenc}
\usepackage[french]{babel}
\usepackage{amsmath,amssymb,amsthm,mathtools}
\usepackage{geometry}
\geometry{margin=2.5cm}
\usepackage{graphicx}
\usepackage{booktabs}
\usepackage{array}
\usepackage{enumitem}
\usepackage{caption}
\usepackage[dvipsnames]{xcolor}
\usepackage{tcolorbox}
\tcbuselibrary{skins,breakable}
\usepackage{hyperref}
\hypersetup{colorlinks=true,linkcolor=MidnightBlue,
            citecolor=MidnightBlue,urlcolor=MidnightBlue}
\usepackage{cleveref}

%--- Environnements théorème --------------------------------------------------
\newtheorem{theoreme}{Théorème}[section]
\newtheorem{proposition}[theoreme]{Proposition}
\newtheorem{lemme}[theoreme]{Lemme}
\newtheorem{corollaire}[theoreme]{Corollaire}
\theoremstyle{definition}
\newtheorem{definition}[theoreme]{Définition}
\newtheorem{probleme}[theoreme]{Problème ouvert}
\theoremstyle{remark}
\newtheorem{remarque}[theoreme]{Remarque}

%--- Boîtes colorées ----------------------------------------------------------
\newtcolorbox{rigorousbox}[1][]{
  colback=Green!8, colframe=Green!60!black, arc=1mm,
  title=\textbf{Résultat rigoureux}, fonttitle=\bfseries, breakable, #1}
\newtcolorbox{resultbox}[1][]{
  colback=RoyalBlue!5, colframe=RoyalBlue!70!black, arc=1mm,
  title=\textbf{Résultat calculatoire}, fonttitle=\bfseries, breakable, #1}
\newtcolorbox{openprobbox}[1][]{
  colback=Goldenrod!10, colframe=Goldenrod!70!black, arc=1mm,
  title=\textbf{Problème ouvert}, fonttitle=\bfseries, breakable, #1}

%--- Raccourcis ---------------------------------------------------------------
\newcommand{\N}{\mathbb{N}}
\newcommand{\Z}{\mathbb{Z}}
\newcommand{\F}{\mathcal{F}}
\newcommand{\Leg}[2]{\left(\dfrac{#1}{#2}\right)}

%==============================================================================
\title{Certificats de primalité inconditionnels\\
pour les entiers de la forme $p = 3m(m+1)+1$ :\\[4pt]
{\large Filtres arithmétiques, recherche multi-c\oe{}urs\\
et records jusqu'à $29\,998$ chiffres décimaux}}

\author{Hassane \textsc{Bakkaoui}\\
\small Chercheur indépendant\\
\small \texttt{bakkahassa@hotmail.com}}

\date{2026}
%==============================================================================

\begin{document}
\maketitle

%------------------------------------------------------------------------------
\begin{abstract}
Nous étudions la sous-famille paramétrique $p = 3m(m+1)+1$,
$m = 2^a\cdot 3^b - 1$, $a,b \in \N^*$, du point de vue de la
certification \emph{inconditionnelle} de primalité via le critère de
Pocklington--Lehmer \cite{Pocklington1914,LehmerPSW1975}.

La structure $3$-lisse de $m+1 = 2^a\cdot 3^b$ fournit, pour tout $(a,b)$,
un diviseur complètement factorisé $F = 2^a\cdot 3^{b+1}$ de $p-1$
satisfaisant $F>\sqrt{p}$ inconditionnellement
(Lemme~\ref{lem:factorisation}), réduisant le certificat à la recherche
de témoins pour $q\in\{2,3\}$ seulement.

Nous établissons trois théorèmes de structure arithmétique :
\begin{itemize}
\item \textbf{Théorème~\ref{thm:mod6} :}
  $p\equiv 1\pmod{6}$ pour tout $(a,b)$.
\item \textbf{Théorème~\ref{thm:cubique} :}
  tout premier $q\equiv 2\pmod{3}$ vérifie $q\nmid p$,
  réduisant le crible de $50\%$.
\item \textbf{Théorème~\ref{thm:mod7} :}
  $7\mid p$ si et seulement si $(a\bmod 3,\,b\bmod 6)\in\F$,
  éliminant $33\%$ des candidats à coût arithmétique nul.
\end{itemize}

Combines, ces trois résultats éliminent $\approx 87\%$ de tous les
candidats avant tout calcul en grande précision. L'algorithme
multi-c\oe{}urs \textsc{PrimeQuest} (versions v1--v7) exploite ce pipeline
pour produire quatre certificats de Pocklington--Lehmer inconditionnels,
le plus grand étant
\[
  p = 3m(m+1)+1,\quad m = 2^{19435}\cdot 3^{19173}-1,
\]
un nombre premier de $\mathbf{29\,998}$ \textbf{chiffres décimaux} trouvé
en $3$ h $06$ min sur un ordinateur portable grand public
(ARM64, 9 c\oe{}urs).

\smallskip
\noindent\textbf{MSC 2020 :} 11A41, 11Y11, 11N13.\\
\textbf{Mots-clés :} certificat de primalité, Pocklington--Lehmer,
filtres arithmétiques, résidus quadratiques, records de grands premiers.
\end{abstract}

\tableofcontents

%==============================================================================
\section{Introduction}\label{sec:intro}
%==============================================================================

Prouver inconditionnellement qu'un entier donné est premier est une tâche
centrale de la théorie algorithmique des nombres. Pour les grands premiers
génériques, la voie la plus pratique est le critère de
Pocklington--Lehmer~\cite{Pocklington1914,LehmerPSW1975} : exhiber un
diviseur complètement factorisé $F>\sqrt{p}$ de $p-1$ et vérifier des
témoins.

La sous-famille
\[
  p = 3m(m+1)+1,\qquad m = 2^a\cdot 3^b-1,\qquad a,b\in\N^*,
\]
est construite de sorte que $m+1=2^a\cdot 3^b$ soit $3$-lisse par
construction, rendant $F=2^a\cdot 3^{b+1}$ complètement factorisé et
vérifiant $F>\sqrt{p}$ pour tout $(a,b)$, sans aucune condition sur $p$.

Au-delà de la structure du certificat, la famille satisfait trois
contraintes arithmétiques indépendantes
(Théorèmes~\ref{thm:mod6}--\ref{thm:mod7}) qui éliminent collectivement
$87\%$ des candidats avant tout calcul en grande précision.

%==============================================================================
\section{La sous-famille de Pocklington}\label{sec:famille}
%==============================================================================

\begin{definition}[Sous-famille de Pocklington]\label{def:pocklington}
Pour des entiers $a,b\geq 1$, on définit
\[
  m(a,b) = 2^a\cdot 3^b - 1, \qquad
  p(a,b) = 3m(m+1)+1 = 3m^2+3m+1.
\]
\end{definition}

\begin{rigorousbox}
\begin{lemme}[Factorisation explicite]\label{lem:factorisation}
$p - 1 = 2^a\cdot 3^{b+1}\cdot m$, $F := 2^a\cdot 3^{b+1} > \sqrt{p}$.
\end{lemme}
\end{rigorousbox}

\begin{rigorousbox}
\begin{theoreme}[Certificat de Pocklington--Lehmer
\cite{Pocklington1914,LehmerPSW1975}]\label{thm:Pocklington}
$p$ est premier ssi $\exists\,w_2,w_3$ tels que, pour $q\in\{2,3\}$,
$w_q^{p-1}\equiv 1\pmod{p}$ et $\gcd(w_q^{(p-1)/q}-1,p)=1$.
\end{theoreme}
\end{rigorousbox}

%==============================================================================
\section{Théorèmes de structure arithmétique}\label{sec:filtres}
%==============================================================================

\begin{rigorousbox}
\begin{theoreme}[Modulo $6$]\label{thm:mod6}
Pour tout $a,b\geq 1$, $p(a,b)\equiv 1\pmod{6}$.
\end{theoreme}
\end{rigorousbox}

\begin{rigorousbox}
\begin{theoreme}[Crible par réciprocité cubique]\label{thm:cubique}
Soit $q\geq 5$ premier avec $q\equiv 2\pmod{3}$. Alors $q\nmid p$.
\end{theoreme}
\end{rigorousbox}

\begin{corollaire}[Réduction du crible de $50\%$]\label{cor:crible-50}
Seuls les premiers $q\equiv 1\pmod{3}$ doivent être testés.
\end{corollaire}

\begin{rigorousbox}
\begin{theoreme}[Classes interdites modulo $7$]\label{thm:mod7}
Soit $\F_7 = \{(0,2),(0,3),(1,0),(1,1),(2,4),(2,5)\}$.
Alors $7\mid p(a,b) \Leftrightarrow (a\bmod 3,\,b\bmod 6)\in\F_7$.
\end{theoreme}
\end{rigorousbox}

%==============================================================================
\section{L'algorithme \textsc{PrimeQuest}}\label{sec:algo}
%==============================================================================

\begin{table}[ht]\centering
\caption{Versions successives de \textsc{PrimeQuest}.}
\label{tab:versions}
\begin{tabular}{@{}clcc@{}}
\toprule
Version & Optimisation ajoutée & Tours MR $k$ & Limite crible $L$ \\
\midrule
v1 & Recherche séquentielle de base & $25$ & $10^6$ \\
v2 & Crible $q\equiv 1\pmod{3}$ & $25$ & $10^6$ \\
v3 & Multi-c\oe{}urs, point de reprise & $25$ & $10^6$ \\
v4 & Filtre mod-$7$ & $25$ & $10^6$ \\
v5 & $k\colon 25\to 3+20$ ; $L\colon 10^6\to 10^7$ & $3+20$ & $10^7$ \\
v6 & Exploration multi-ratio (round-robin) & $3+20$ & $10^7$ \\
v7 & Tolérance $\pm 50$, spread adaptatif & $3+20$ & $10^7$ \\
\bottomrule
\end{tabular}
\end{table}

%==============================================================================
\section{Résultats calculatoires}\label{sec:resultats}
%==============================================================================

\begin{table}[ht]\centering
\caption{Certificats Pocklington de la campagne \textsc{PrimeQuest}.}
\label{tab:resultats}
\begin{tabular}{@{}rrrrrrrr@{}}
\toprule
Chiffres de $p$ & $a$ & $b$ & $b/a$ & Tests MR & Élim. & Temps & Workers \\
\midrule
$9\,998$  & $6\,212$  & $6\,738$  & $1{,}085$ & ---      & ---      & ---             & $1$ \\
$10\,000$ & $6\,213$  & $6\,740$  & $1{,}085$ & ---      & ---      & ---             & $1$ \\
$19\,999$ & $12\,228$ & $13\,242$ & $1{,}083$ & $1\,050$ & $87{,}9\%$ & $2$\,h\,$40$\,min & $9$ \\
$29\,998$ & $19\,435$ & $19\,173$ & $0{,}987$ & $286$    & $87{,}1\%$ & $3$\,h\,$06$\,min & $9$ \\
\bottomrule
\end{tabular}
\end{table}

%==============================================================================
\section{Conclusion}\label{sec:conclusion}
%==============================================================================

Trois théorèmes arithmétiques élémentaires éliminent $87\%$ des candidats
avant tout calcul en grande précision. Appliqué via l'algorithme
multi-c\oe{}urs \textsc{PrimeQuest}, ce cadre a produit quatre certificats
inconditionnels, culminant en un premier de $\mathbf{29\,998}$ chiffres
décimaux. Une cinquième recherche ciblant $50\,000$ chiffres est en cours.

%==============================================================================
\begin{thebibliography}{99}
%==============================================================================

\bibitem{LehmerPSW1975}
J.~Brillhart, D.~H.~Lehmer et J.~L.~Selfridge,
Math.\ Comp.\ \textbf{29} (1975), 620--647.

\bibitem{Pocklington1914}
H.~C.~Pocklington,
Proc.\ Cambridge Philos.\ Soc.\ \textbf{18} (1914--1916), 29--30.

\bibitem{IrelandRosen}
K.~Ireland et M.~Rosen, \emph{A Classical Introduction to Modern
Number Theory}, Springer, 1990.

\end{thebibliography}

\end{document}
